\section{Introduction}
A scene query specifies both an identity and the spatial extent to be returned.
A text description may match several observations, a reference image may expose
only part of an object, and a point identifies one surface rather than every
member. Modern visual teachers provide strong local correspondences. Turning
these correspondences into a complete three-dimensional answer still requires
deciding what each observation actually establishes.

A reliable match and a complete observation are different properties. A mask
may correctly identify a flower while omitting its leaves. Transporting that
mask accurately can reproduce the omission across views; treating its exterior
as background can actively suppress the missing members. Conversely, expanding
the response without preserving identity can capture a neighboring object.
The scientific question is therefore: \emph{when does partial query evidence
justify support or exclusion, and how should this evidence form a complete
response under the intended output semantics?}

We address this question with \emph{query-conditioned evidence fusion} (QCEF).
The central mechanism assigns evidence authority before aggregation: identity
applicability determines which observations can act, and spatial scope
determines where they can support, exclude or leave the response unknown.
Native teacher scores remain in their own matching spaces. Admitted responses
then enter a geometric fusion or conditional-membership operator appropriate
to the requested output. This is a shared evidence-to-response formulation
with typed realizations, not a claim that every modality uses one network or
one inverse solver.

Reference queries make the mechanism particularly concrete. We complete the
reference response using signed query evidence, select between whole reference
candidates using reference-side feasibility or predicted quality, and only then
transport the selected response to other views. This changes the scope of the
transmitted evidence rather than asking a stronger reconstruction objective to
repair incomplete targets. For map-only image and point queries, the analogous
extent requirement is implemented by conditional membership over the complete
carrier. Complete here describes the output domain and intended capability;
it does not imply that every predicted object is recovered without error.

The analysis distinguishes three obstacles. Partial labels alone need not
identify complete extent. Incorrect negative authority attenuates a weighted
response even when correspondences are correct. Finally, rendering a selected
set and averaging membership in the full scene can disagree under occlusion.
These statements motivate scope estimation and typed objectives, without
turning fallible teacher predictions into guaranteed semantic truth.

Our implementation includes the established system recorded as v10 and the
reference-completion extension. Its evidence spans LERF semantic segmentation
and projected selection, ScanNet open-vocabulary segmentation, NVOS reference
segmentation and multimodal UQIS. Historical controls isolate identity,
response-use and membership capabilities. Matched reference interventions show
that completing scope before transport is useful even without target-image
evidence; stronger source fitting alone has mixed effects. Published results
and verified local reproductions provide comparison context with their input
and metric differences retained.

Our contributions are:
\begin{enumerate}
\item A response-level formulation of partial-evidence authority, with scoped
analysis of extent ambiguity, negative-evidence attenuation and output
non-equivalence. It separates the information required to identify an object
from the evidence needed to recover its members.
\item A query-conditioned fusion method that retains native matching while
combining admitted evidence through typed geometric and membership operators.
Reference completion and whole-response admission before transport provide a
concrete mechanism for repairing incomplete query scope.
\item Five-output evaluation with matched component interventions, protocol
audits and counterexamples to stronger-fitting explanations. The evidence
connects measured gains to specific capabilities and records where the shared
formulation has not yet yielded a universally effective admission rule.
\end{enumerate}
